\documentclass[11pt]{article}
\usepackage[letterpaper,margin=0.83in,headheight=16pt]{geometry}
\usepackage[T1]{fontenc}
\usepackage[utf8]{inputenc}
\usepackage{lmodern}
\usepackage{microtype}
\usepackage{amsmath,amssymb}
\usepackage{xcolor}
\usepackage{hyperref}
\usepackage{enumitem}
\usepackage{booktabs,longtable,tabularx,array}
\usepackage{fancyhdr}
\usepackage{titlesec}
\usepackage{listings}
\usepackage{tikz}
\usepackage{float}
\definecolor{suiteblue}{HTML}{17365D}
\definecolor{suitelight}{HTML}{EAF1F8}
\definecolor{suitegray}{HTML}{556070}
\hypersetup{colorlinks=true,linkcolor=suiteblue,urlcolor=suiteblue,breaklinks=true}
\titleformat{\section}{\Large\bfseries\color{suiteblue}}{\thesection}{0.55em}{}
\titleformat{\subsection}{\large\bfseries\color{suiteblue}}{\thesubsection}{0.55em}{}
\setlist[itemize]{itemsep=0.25em,topsep=0.35em,leftmargin=1.4em}
\setlist[enumerate]{itemsep=0.25em,topsep=0.35em,leftmargin=1.6em}
\setlength{\parindent}{0pt}
\setlength{\parskip}{0.47em}
\emergencystretch=2em
\lstset{basicstyle=\ttfamily\footnotesize,breaklines=true,frame=single,columns=fullflexible,showstringspaces=false}
\pagestyle{fancy}\fancyhf{}
\lhead{Game Design \& Engineering}\rhead{Playtesting Procedure}\cfoot{\thepage}
\begin{document}
{\huge\bfseries\color{suiteblue} Playtest Execution And Observation Procedure}\par
\vspace{0.15in}
{\normalsize\color{suitegray} Original operational process; Schell Chapter 28 is the conceptual source}\par
\vspace{0.15in}\hrule\vspace{0.15in}

\section{Purpose and source boundary}
Book-aligned topic: the six core questions of playtesting in Chapter 28. Consent mechanics, measurements, ownership, classification, and evidence handling are supplementary implementation guidance.


\section{Step 1: Environment}
Record build commit, runtime, venue, controls, accessibility settings, input hardware, and environmental interruptions.

Exit evidence: dated, attributable process record and a verification link or explicit reason for omission.


\section{Step 2: Introduction}
Use a consistent neutral script. Explain withdrawal, data handling, and the permitted level of assistance.

Exit evidence: dated, attributable process record and a verification link or explicit reason for omission.


\section{Step 3: Observations}
Timestamp action, context, outcome, prompt, and exact player quote only when authorized. Keep interpretation in a separate column.

Exit evidence: dated, attributable process record and a verification link or explicit reason for omission.


\section{Step 4: Intervention}
Prompt only under a predeclared facilitation rule; a prompted success is not the same as unassisted success.

Exit evidence: dated, attributable process record and a verification link or explicit reason for omission.


\section{Step 5: Debrief}
Ask open questions about goals, surprises, interpretation of mechanics, frustrations, and unexpected strategies.

Exit evidence: dated, attributable process record and a verification link or explicit reason for omission.


\section{Step 6: Clean-up}
Close session, confirm consent status, pseudonymize materials, and set deletion deadline.

Exit evidence: dated, attributable process record and a verification link or explicit reason for omission.


\section{Required controls}
\begin{itemize}
\item Reproducible build and hardware record
\item Valid opt-in consent and participant withdrawal mechanism
\item Event log separating observation, quote, inference, and recommendation
\item Documented safety stop and accessibility accommodation procedure
\item Finding priority register with reviewer and re-test date
\end{itemize}

\section{Process scorecard}
Score each control as Complete, Partial, Not Demonstrated, or Not Applicable (with rationale). Do not convert absent participant evidence into an automatic pass.


\end{document}
